% Figure: Newtonian-style reflecting telescope layout. A parallel bundle from a
% distant star strikes the concave primary mirror of diameter D and focal length
% f, converges toward the prime focus, is intercepted by a flat diagonal
% secondary, and is brought to the focal plane at the eyepiece. The drawing fixes
% the focal ratio f/D, the converging cone, and the position of the detector.
\begin{figure}[htbp]
\centering
\begin{tikzpicture}[scale=1.0]
% optical axis
\draw[enggray, dashed] (-1,0) -- (11,0);
% primary mirror (concave) at right, drawn as a vertical bracket
\draw[engblue, very thick] (10,-2.2) .. controls (9.4,-1.1) and (9.4,1.1) .. (10,2.2);
\node[engblue, font=\scriptsize, anchor=west] at (10.1,0) {primary $D$};
% incoming parallel rays from a distant star
\foreach \y in {-1.8,-1.2,-0.6,0.6,1.2,1.8} {
  \draw[engorange, thick] (-1,\y) -- (9.55,\y);
}
% converging cone from primary toward prime focus near x=3
\draw[engorange, thick] (9.55,1.8) -- (3.6,0.15);
\draw[engorange, thick] (9.55,1.2) -- (3.6,0.10);
\draw[engorange, thick] (9.55,0.6) -- (3.6,0.05);
\draw[engorange, thick] (9.55,-0.6) -- (3.6,-0.05);
\draw[engorange, thick] (9.55,-1.2) -- (3.6,-0.10);
\draw[engorange, thick] (9.55,-1.8) -- (3.6,-0.15);
% flat diagonal secondary intercepts cone at x about 4.2
\draw[engred, very thick] (4.6,-0.7) -- (3.8,0.7);
\node[engred, font=\scriptsize, anchor=south] at (4.2,0.75) {flat secondary};
% folded rays to side focal plane
\draw[engorange, thick] (4.25,0.32) -- (4.25,2.0);
\draw[engorange, thick] (4.15,-0.08) -- (4.15,2.0);
% focal plane / detector
\draw[engblack, very thick] (3.7,2.0) -- (4.7,2.0);
\node[engblack, font=\scriptsize, anchor=south] at (4.2,2.05) {focal plane};
% focal length brace along axis
\draw[englime!60!engblack, |<->|, thick] (3.6,-2.0) -- (9.55,-2.0);
\node[englime!60!engblack, font=\scriptsize, anchor=north] at (6.5,-2.05)
  {focal length $f$};
\node[engorange, font=\scriptsize, anchor=south east] at (-0.1,1.8)
  {from a distant star};
\end{tikzpicture}
\caption[Reflecting-telescope layout]{A reflecting telescope folds a parallel
star bundle to a focal plane. The concave primary mirror of clear aperture $D$
and focal length $f$ gathers light over an area $\pi D^2/4$ and converges it
toward the prime focus; a small flat secondary intercepts the converging cone
and folds it to an accessible focal plane where the detector sits. Two numbers
set the design: the aperture $D$, which fixes both the light-gathering area and
the diffraction limit, and the focal ratio $f/D$, which fixes the plate scale
and the speed. The secondary obscures a small central fraction of the aperture,
a loss the area budget must carry.}
\label{fig:vol04ch11:telescope-layout}
\end{figure}
